\documentclass[11pt]{article}

\usepackage{amsmath,amssymb}
\usepackage[a4paper,margin=1in]{geometry}

% Remove paragraph indentation
\setlength{\parindent}{0pt}

\begin{document}

\title{A Concise Derivation of Principal Component Analysis}
\author{Zhenwei Tang}
\date{}
\maketitle

\section*{Goal}

Given data points
\[
x_i \in \mathbb{R}^p,
\qquad i=1,\dots,n,
\]
PCA aims to construct a lower-dimensional representation
\[
y_i = U_k^T(x_i-m),
\qquad
y_i \in \mathbb{R}^k,
\qquad
U_k \in \mathbb{R}^{p\times k},
\qquad
k < p.
\]

Here, \(m\) is the mean of the data, and the columns of \(U_k\) define the \(k\)-dimensional subspace onto which the data is projected.

% some more text here ...

\section*{Step 1: Center the Data}

The mean of the data is
\[
m = \frac{1}{n}\sum_{i=1}^n x_i.
\]

% some more text here ...

\section*{Conclusion}

PCA starts from the goal
\[
y_i = U_k^T(x_i-m),
\]
where \(y_i\) is a lower-dimensional representation of \(x_i\).

To preserve as much information as possible, PCA maximizes the variance of the projected data. This leads to the constrained optimization problem
\[
\max_v v^TSv
\quad
\text{subject to}
\quad
v^Tv=1.
\]

% some more text here ...

\[
\boxed{
\text{PCA reduces dimensionality by projecting data onto the directions of maximal variance.}
}
\]

\end{document}